\begin{lemma}
\label{lemma-homotopy}
Let $f^0, f^1 : V \to U$ be maps of simplicial sets.
Let $n \geq 0$ be an integer.
Assume
\begin{enumerate}
\item The maps $f^j_i : V_i \to U_i$, $j = 0, 1$ are equal for $i < n$.
\item The canonical morphism $U \to \text{cosk}_n \text{sk}_n U$
is an isomorphism.
\item The canonical morphism $V \to \text{cosk}_n \text{sk}_n V$
is an isomorphism.
\end{enumerate}
Then $f^0$ is homotopic to $f^1$.
\end{lemma}

\begin{proof}[First proof]
Let $W$ be the $n$-truncated simplicial set with $W_i = U_i$ for
$i < n$ and $W_n = U_n / \sim$ where $\sim$ is the equivalence relation
generated by $f^0(y) \sim f^1(y)$ for $y \in V_n$. This makes sense
as the morphisms $U(\varphi) : U_n \to U_i$ corresponding to
$\varphi : [i] \to [n]$ for $i < n$ factor through the quotient map
$U_n \to W_n$ because $f^0$ and $f^1$ are morphisms of simplicial sets and
equal in degrees $< n$. Next, we upgrade $W$ to a simplicial set
by taking $\text{cosk}_n W$. By Lemma \ref{lemma-section}
the morphism $g : U \to W$ is a trivial Kan fibration. Observe
that $g \circ f^0 = g \circ f^1$ by construction and denote
this morphism $f : V \to W$. Consider the diagram
$$
\xymatrix{
\partial \Delta[1] \times V \ar[rr]_{f^0, f^1} \ar[d] & & U \ar[d] \\
\Delta[1] \times V \ar[rr]^f \ar@{-->}[rru] & & W
}
$$
By Lemma \ref{lemma-trivial-kan} the dotted arrow exists and the proof is done.
\end{proof}

\begin{proof}[Second proof]
We have to construct a morphism of simplicial sets
$h : V \times \Delta[1] \to U$ which recovers $f^i$ on composing with $e_i$.
The case $n = 0$ was dealt with above the lemma.
Thus we may assume that $n \geq 1$.
The map $\Delta[1] \to \text{cosk}_1 \text{sk}_1 \Delta[1]$
is an isomorphism, see Lemma \ref{lemma-simplex-cosk}.
Thus we see that $\Delta[1] \to \text{cosk}_n \text{sk}_n \Delta[1]$
is an isomorphism as $n \geq 1$, see
Lemma \ref{lemma-cosk-up}. And hence $V \times \Delta[1] \to
\text{cosk}_n \text{sk}_n (V \times \Delta[1])$
is an isomorphism too, see Lemma \ref{lemma-cosk-product}.
In other words, in order to construct the homotopy
it suffices to construct a suitable
morphism of $n$-truncated simplicial sets
$h : \text{sk}_n V \times \text{sk}_n \Delta[1] \to \text{sk}_n U$.

\medskip\noindent
For $k = 0, \ldots, n - 1$ we define $h_k$ by the
formula $h_k(v, \alpha) = f^0(v) = f^1(v)$.
The map $h_n : V_n \times \Mor_{\Delta}([n], [1]) \to U_n$
is defined as follows. Pick $v \in V_n$ and $\alpha : [n] \to [1]$:
\begin{enumerate}
\item If $\Im(\alpha) = \{0\}$, then we set $h_n(v, \alpha) = f^0(v)$.
\item If $\Im(\alpha) = \{0, 1\}$, then we set $h_n(v, \alpha) = f^0(v)$.
\item If $\Im(\alpha) = \{1\}$, then we set $h_n(v, \alpha) = f^1(v)$.
\end{enumerate}
Let $\varphi : [k] \to [l]$ be a morphism of $\Delta_{\leq n}$.
We will show that the diagram
$$
\xymatrix{
V_{l} \times \Mor([l], [1]) \ar[r] \ar[d] & U_{l} \ar[d] \\
V_{k} \times \Mor([k], [1]) \ar[r] & U_{k}
}
$$
commutes.
Pick $v \in V_{l}$ and $\alpha : [l] \to [1]$.
The commutativity means that
$$
h_k(V(\varphi)(v), \alpha \circ \varphi)
=
U(\varphi)(h_l(v, \alpha)).
$$
In almost every case this holds because
$h_k(V(\varphi)(v), \alpha \circ \varphi) = f^0(V(\varphi)(v))$
and $U(\varphi)(h_l(v, \alpha)) = U(\varphi)(f^0(v))$, combined
with the fact that $f^0$ is a morphism of simplicial sets.
The only cases where this does not hold is when
either (A) $\Im(\alpha) = \{1\}$ and $l = n$
or (B) $\Im(\alpha \circ \varphi) = \{1\}$ and $k = n$.
Observe moreover that necessarily $f^0(v) = f^1(v)$
for any degenerate $n$-simplex of $V$.
Thus we can narrow the cases above down even further
to the cases (A) $\Im(\alpha) = \{1\}$, $l = n$
and $v$ nondegenerate, and (B)
$\Im(\alpha \circ \varphi) = \{1\}$, $k = n$
and $V(\varphi)(v)$ nondegenerate.

\medskip\noindent
In case (A), we see that also $\Im(\alpha \circ \varphi) = \{1\}$.
Hence we see that not only $h_l(v, \alpha) = f^1(v)$ but also
$h_k(V(\varphi)(v), \alpha \circ \varphi) = f^1(V(\varphi)(v))$.
Thus we see that the relation holds because $f^1$ is a morphism
of simplicial sets.

\medskip\noindent
In case (B) we conclude that $l = k = n$ and
$\varphi$ is bijective, since otherwise $V(\varphi)(v)$
is degenerate. Thus $\varphi = \text{id}_{[n]}$, which is a trivial case.
\end{proof}